\subsection*{集合}
\begin{enumerate}
    \item 设集合$A = \left\{1, 4, x\right\}$，$B = \left\{1, x^2\right\}$，且$A \cap B = B$，则$x = $\blkda[3]{$0 \mbox{或} -2 \mbox{或} 2$}
    \begin{jieda}
        首先对条件进行转换，$A \cap B = B$即$B \subseteq A$，根据集合中元素的无序性可知需要进行分类讨论。
        \begin{enumerate}
            \item $x^2 = x$，此时对应于$x = 0 \mbox{或} 1$，但是经过检验发现$x = 1$使得集合中元素违背互异性。
            \item $x^2 = 4$，此时对应于$x = 2 \mbox{或} -2$，均不违背互异性。
        \end{enumerate}
        综上，$x = 0 \mbox{或} -2 \mbox{或} 2$
    \end{jieda}
    \item 若集合$S = \left\{x \left| \dfrac{6}{x-2}\ \in \mathbf{Z} \mbox{且} x \in \mathbf{Z}\right. \right\}$, 则集合S的非空真子集个数为\blkda{$254$}.
    \begin{jieda}
        由于 $\dfrac{6}{x-2}\ \in \mathbf{Z}$, 因此$x-2 \in \left\{-6, -3, -2, -1, 1, 2, 3, 6\right\}$，即 $S = \left\{-4, -1, 0, 1, 3, 4, 5, 8\right\}$， 因此集合$S$的元素个数为8， 因此其非空真子集的个数为$2^8-2=254$ \\
        \ps 含有n个元素的有限集的非空真子集个数为$2^n -2$
    \end{jieda}

    \item 已知集合 $A = \left\{  {x\left| {a{x}^{2} - {3x} + 2 = 0}\right. }\right\}$ 中有且仅有一个元素，则 $a$ 的取值组成的集合为\blkda{$\left\{0, \dfrac{9}{8}\right\}$}.
    \begin{jieda}
        $a = 0$或$\Delta = 0$
    \end{jieda}

    \item 若规定集合 $E = \{ 0,1,2,\cdots ,n\}$ 的子集 $\left\{  {{a}_{1},{a}_{2},{a}_{3},\cdots ,{a}_{m}}\right\}$ 为 $E$ 的第 $k$ 个子集，其中 $k = {2}^{{a}_{1}} + {2}^{{a}_{2}} + {2}^{{a}_{3}} + \cdots  + {2}^{{a}_{m}}$ ，则 $E$ 的第 211 个子集是\blkda[2]{$\{0,1,4,6,7\}$}.
    \begin{jieda}
        $211 = 128 + 64 + 16 + 2 + 1$
    \end{jieda}
    \item 设集合 $A = \{ x \mid   - 2 \leq  x \leq  5\} ,B = \{ x \mid  m + 1 \leq  x \leq  {2m} - 1\}$.
    \begin{enumerate}
        \item 若 $A \cap  B = B$ ，求实数 $m$ 的取值范围
        \item 当 $x \in  \mathbf{R}$ 时，不存在 $x$ 使得 $x \in  A$ 与 $x \in  B$ 同时成立，求实数 $m$ 的取值范围.
    \end{enumerate}
    \begin{jieda}
        \begin{enumerate}
            \item $\because A \cap  B = B,\therefore B \subseteq  A$
            
            当 $m + 1 > {2m} - 1$ ,即 $m < 2$ 时, $B = \varnothing$ ,满足 $B \subseteq  A$
            
            当 $m + 1 \leq  {2m} - 1$ ,即 $m \geq  2$ 时,要使 $B \subseteq  A$ 成立,必须且只需 $\left\{  \begin{array}{l} m + 1 \geq   - 2, \\  {2m} - 1 \leq  5, \end{array}\right.$ 可得 $2 \leq  m \leq  3$ ,综上,若 $A \cap  B = B$ ,则 $m \leq  3$ .
            \item 依题意, $A \cap  B = \varnothing$ . 
            
            若 $B = \varnothing$ ,即 $m + 1 > {2m} - 1$ ,得 $m < 2$ 时满足条件；
            
            若 $B \neq  \varnothing$ ,则 $\left\{  \begin{array}{l} m + 1 \leq  {2m} - 1, \\  m + 1 > 5 \end{array}\right.$ 或 $\begin{array}{l} m + 1 \leq  {2m} - 1, \\  {2m} - 1 <  - 2, \end{array}$ 解得 $m > 4$. 
            
            综上,有 $m < 2$ 或 $m > 4$.
        \end{enumerate}
    \end{jieda}
    
    \item 已知 $a,m \in  \mathbf{R}$ ,集合 $A = \left\{  {x \mid  {x}^{2} - {3x} + 2 = 0}\right\} $，$B = \left\{  {x \mid  {x}^{2} + 2\left( {m + 1}\right) x + {m}^{2} - 5 = 0,x \in  \mathbf{R}}\right\} $，$C =  \left\{  {x \mid  {x}^{2} - {ax} + a - 1 = 0,x \in  \mathbf{R}}\right\} .$
    \begin{enumerate}
        \item 若 $A \cup  C = A$ ，求实数 $a$ 的值
        \item 若 $A \cap  B = B$ ，求实数 $m$ 的取值范围
    \end{enumerate}
    \begin{jieda}
        \begin{enumerate}
            \item $A = \{ 1,2\}$，由 ${x}^{2} - {ax} + a - 1 = 0$ 得 $x = 1$ 或 $a - 1$ ,所以 $1 \in  C,a - 1 \in  C$
            
            因为 $A \cup  C = A$ ,所以 $C \subseteq  A$ ,所以 $a - 1 = 1$ 或 2,所以 $a = 2$ 或 3;
            \item 因为 $A \cap  B = B$ ,所以 $B \subseteq  A$
            
            当 $B = \varnothing$ 时, $\Delta  = 4{\left( m + 1\right) }^{2} - 4\left( {{m}^{2} - 5}\right)  < 0$ ,解得 $m <  - 3$
            
            当 $B = \{ 1\}$ 时, $\left\{  \begin{array}{l} \Delta  = 4{\left( m + 1\right) }^{2} - 4\left( {{m}^{2} - 5}\right)  = 0 \\  1 + 2\left( {m + 1}\right)  + {m}^{2} - 5 = 0 \end{array}\right.$ ,无解
            
            当 $B = \{ 2\}$ 时, $\left\{  \begin{array}{l} \Delta  = 4{\left( m + 1\right) }^{2} - 4\left( {{m}^{2} - 5}\right)  = 0 \\  4 + 4\left( {m + 1}\right)  + {m}^{2} - 5 = 0 \end{array}\right.$ ,解得 $m =  - 3$
            
            当 $B = \{ 1,2\}$ 时, $\left\{  \begin{array}{l} 1 + 2 =  - 2\left( {m + 1}\right) \\  1 \cdot  2 = {m}^{2} - 5 \end{array}\right.$ ,无解.
            
            综上,实数 $m$ 的取值范围是 $( - \infty , - 3\rbrack$.
        \end{enumerate}
    \end{jieda}
    
    \item 设 $f\left( x\right)  = {x}^{2} + {ax} + b\left( {a,b \in  \mathbf{R}}\right) ,A = \{ x \mid  f\left( x\right)  = x\} ,B = \left\{  {x\left| {f\left\lbrack  {f\left( x\right) }\right\rbrack   = x}\right. }\right\}$，若 $A = \{  - 1,3\}$ ，求集合 $B$.
    \begin{jieda}
        $A = \{  - 1,3\}  \Longrightarrow  a - 1 =  - 2,b =  - 3 \Longrightarrow  f\left( x\right)  = {x}^{2} - x - 3$，于是 $f\left\lbrack  {f\left( x\right) }\right\rbrack   = {\left( {x}^{2} - x - 3\right) }^{2} - \left( {{x}^{2} - x - 3}\right)  - 3 = x,\therefore B = \{  - 1,3, - \sqrt{3},\sqrt{3}\}$
    \end{jieda}
    
    \item 设 $a$ 、 $b$ 、 $c \in  \mathbf{R}$ ,函数 $f\left( x\right)  = \left( {x + a}\right) \left( {{x}^{2} + {bx} + c}\right) ,g\left( x\right)  = \left( {{ax} + 1}\right) \left( {c{x}^{2} + {bx} + 1}\right)$ . 记集合 $S = \{ x \mid  f\left( x\right)  = 0$ , $x \in  \mathbf{R}\} ,T = \{ x \mid  g\left( x\right)  = 0,x \in  \mathbf{R}\}$ . 若 $\left| S\right| ,\left| T\right|$ 分别为集合 $S,T$ 的元素个数,则下列结论不可能的是\dotfill{(\kh{D})}
    \xx{$\left| S\right|  = 1,\left| T\right|  = 0$}{$\left| S\right|  = 1,\left| T\right|  = 1$}{$\left| S\right|  = 2,\left| T\right|  = 2$}{$\left| S\right|  = 2,\left| T\right|  = 3$}
    \begin{jieda}
        $|S|$最多3个元素，$|T|$最多和$|S|$相等，且元素间有一一对应的倒数关系. 当$a,c$中至少有一个为0， $|T|$比$|S|$少一个元素.
    \end{jieda}

\end{enumerate}
\clearpage
\subsection*{逻辑}
\begin{enumerate}
    \item 设 $a$ 、 $b \in  \mathbf{R}$ ，则“$a + b > 0$”是“$a > 0$ 且 $b > 0$”的 \dotfill{(\kh{B})}
    \xx{充分非必要条件}{必要非充分条件}{充要条件}{既非充分又非必要条件}
    \item 已知$p: x-3 \leq 2, q: (x-m+1)(x-m-1) \leq 0$，若$\neg p$是$\neg q$的充分非必要条件，求实数$m$的取值范围.
    \begin{jieda}
        $p$所对应的集合$P = (-\infty, 5]$，故$\neg p$对应的集合$\overline{P} = (5, +\infty)$ \\
        $q$所对应的集合$Q = [m-1, m+1]$，故$\neq q$对应的集合$\overline{Q} = (-\infty, m-1) \cup (m+1, +\infty)$ \\
        $\neg p$是$\neg q$的充分非必要条件即$\overline{P} \subset \overline{Q}$，即$m+1 \leq 5$，解得$m \in (-\infty, 4]$.
    \end{jieda}

    \item 已知 $A = \{ x \mid  \left| {x - a}\right|  < 4\} ,B = \{ x \mid  \left( {x - 1}\right) \left( {2 - x}\right)  > 0\}$ ,若 $x \in  A$ 是 $x \in  B$ 的必要非充分条件,则实数 $a$ 的取值范围是\blkda{$[-2, 5]$}
    \begin{jieda}
        $A = (a-4, a+4)$，$B = (1,2)$，$B \subset A$
    \end{jieda}
    
    \item 已知集合 $A = \left\{  {y\left| {y = {x}^{2} - \dfrac{3}{2}x + 1}\right. ,x \in  \left\lbrack  {\dfrac{3}{4},2}\right\rbrack  }\right\}  ,B = \left\{  {x\left| {x + {m}^{2} \geq  1}\right. }\right\}$，若 $p : x \in  A;q : x \in  B$ 且 $p$ 是 $q$ 的充分条件，求实数 $m$ 的取值范围.
    \begin{jieda}
        \begin{enumerate}
            \item $y = {x}^{2} - \dfrac{3}{2}x + 1 = {\left( x - \dfrac{3}{4}\right) }^{2} + \dfrac{7}{16},\because x \in  \left\lbrack  {\dfrac{3}{4},2}\right\rbrack  ,\therefore \dfrac{7}{16} \leq  y \leq  2$ ,即集合 $A = \left\lbrack  {\dfrac{7}{16},2}\right\rbrack$ ；
            \item $B = \left\{  {x \mid  x + {m}^{2} \geq  1}\right\}   = \left\{  {x \mid  x \geq  1 - {m}^{2}}\right\}$.
            
            若 $p : x \in  A;q : x \in  B$ 且 $p$ 是 $q$ 的充分条件,则 $A \subseteq  B,1 - {m}^{2} \leq  \dfrac{7}{16}$ ,解得 $m \geq  \dfrac{3}{4}$ 或 $m \leq   - \dfrac{3}{4}$，即实数 $m$ 的取值范围是 $\left\{ m \left|  m \geq  \dfrac{3}{4} \text{或} m \leq   - \dfrac{3}{4} \right. \right\}$.
        \end{enumerate}
    \end{jieda}
    
    \item 已知直线 ${l}_{1} : {ax} + {2y} - a + 1 = 0,{l}_{2} : {3x} + \left( {a - 1}\right) y - 2 = 0$ ,则 “ $a =  - 2$ ” 是 “ ${l}_{1}//{l}_{2}$ ” 的\dotfill{(\kh{C})}
    \xx{充分不必要条件}{必要不充分条件}{充要条件}{既不充分也不必要条件}
    \begin{jieda}
        当 $a =  - 2$ ,直线 ${l}_{1} : {2x} - {2y} - 3 = 0,{l}_{2} : {3x} - {3y} - 2 = 0$ ,此时 ${l}_{1}//{l}_{2}$ ,故 “ $a =  - 2$ ” 是 “ ${l}_{1}//{l}_{2}$ ” 的充分条件

        由 ${l}_{1}//{l}_{2}$ ,得 $\left\{  \begin{array}{l} a\left( {a - 1}\right)  - 2 \times  3 = 0 \\   - {2a} - 3\left( {1 - a}\right)  \neq  0 \end{array}\right.$ ,解得 $a =  - 2$ ,故 “ $a =  - 2$ ” 是 “ ${l}_{1}//{l}_{2}$ ” 的必要条件

        故 “ $a =  - 2$ ” 是 “ ${l}_{1}//{l}_{2}$ ” 的充要条件.
    \end{jieda}

    \item 在 $\bigtriangleup  {ABC}$ 中， $P$ 在线段 ${BC}$ 上，满足 $2\vv{BP} = \vv{PC}$ ， $O$ 是线段 ${AP}$ 的中点. 过点 $O$ 的直线与线段 ${AB},{AC}$ 分别交于点 $E,F$，设 $\vv{AE} = \lambda \vv{AB}$ ， $\vv{AF} = \mu \vv{AC}$.

    \begin{center}
    \includegraphics[width=0.2\textwidth]{bodyxb/src/1_150_1623_342_219_0.jpg}
    \end{center}
    判断以下命题中正确的是\dotfill{(\kh{C})}
    \begin{dingautolist}{172}
        \item $\dfrac{2}{\lambda } + \dfrac{1}{\mu }$ 为定值
        \item 设 $\bigtriangleup {AEF}$ 的面积为 ${S}_{1},\bigtriangleup {ABC}$ 的面积为 ${S}_{2}$ ,求 $\dfrac{{S}_{1}}{{S}_{2}}$ 的取值范围是$\left[\dfrac{1}{4}, \dfrac{2}{5}\right]$
    \end{dingautolist}
    \xx{命题 \ding{172} 不正确，命题 \ding{173} 正确}{命题 \ding{172} 、命题 \ding{173} 都不正确}{命题 \ding{172} 正确，命题 \ding{173} 不正确}{命题 \ding{172} 、命题 \ding{173} 都正确}
    \begin{jieda}
        由 $2\vv{BP} = \vv{PC}$ ,得 $\vv{AP} = \vv{AB} + \vv{BP} = \vv{AB} + \dfrac{1}{3}\vv{BC} = \vv{AB} + \dfrac{1}{3}\left( {\vv{BA} + \vv{AC}}\right)  = \dfrac{2}{3}\vv{AB} + \dfrac{1}{3}\vv{AC} = \dfrac{2}{3\lambda }\vv{AE} + \dfrac{1}{3\mu }\vv{AF}$

        由 $O$ 是线段 ${AP}$ 的中点,得 $\vv{AO} = \dfrac{1}{2}\vv{AP} = \dfrac{1}{3\lambda }\vv{AE} + \dfrac{1}{6\mu }\vv{AF}$，由 $E,O,F$ 三点共线,得 $\dfrac{1}{3\lambda } + \dfrac{1}{6\mu } = 1$ ,则 $\dfrac{2}{\lambda } + \dfrac{1}{\mu } = 6$

        ${S}_{1} = \dfrac{1}{2}\left| {AE}\right| \left| {AF}\right| \sin \angle {EAF} = \dfrac{1}{2}\lambda \left| {AB}\right|  \cdot  \mu \left| {AC}\right| \sin \angle {BAC} = {\lambda \mu }{S}_{2}$ ,得 $\dfrac{{S}_{1}}{{S}_{2}} = {\lambda \mu }$

        $\lambda  > 0,\mu  > 0,6 = \dfrac{2}{\lambda } + \dfrac{1}{\mu } \geq  2\sqrt{\dfrac{2}{\lambda } \cdot  \dfrac{1}{\mu }}$ ,即 ${\lambda \mu } \geq  \dfrac{2}{9}$ ,当且仅当 $\dfrac{2}{\lambda } = \dfrac{1}{\mu } = 3$ 时取等号
    \end{jieda}

    \item 2021 年 3 月 30 日，小米正式开始启用具备“超椭圆”数学之美的 LOGO(如图所示)，设计师的灵感来源于曲线 $C : {\left| \dfrac{x}{a}\right| }^{n} + {\left| \dfrac{y}{b}\right| }^{n} = 1\left( {n > 0,n \in  \mathbf{R}}\right)$ . 当 $n = 4,a = 3,b = 2$ 时,给出下列四个结论:

    \begin{center}
    
\includegraphics[width=0.1\textwidth]{bodyxb/src/0_163_428_318_317_0.jpg}
    \end{center}
    \begin{dingautolist}{172}
    \item 曲线 $C$ 关于原点对称
    \item 曲线 $C$ 所围成的封闭图形的面积小于 24
    \item 曲线 $C$ 上的点到原点 $O$ 的距离的最大值为 3
    \item 设 $M\left( {\sqrt{5},0}\right)$ ,直线 $x - y + \sqrt{5} = 0$ 交曲线 $C$ 于 $P,Q$ 两点,则 $\bigtriangleup {PQM}$ 的周长大于 12.
    \end{dingautolist}
    其中所有正确的结论有\blkda{\ding{172}\ding{173}\ding{175}}(填序号).
    \begin{jieda}
    \begin{dingautolist}{172}
    \item 取曲线上点 $\left( {x,y}\right)$ ，则 $\left( {-x, - y}\right)$ 满足 $\dfrac{{\left( -x\right) }^{4}}{81} + \dfrac{{\left( -y\right) }^{4}}{16} = \dfrac{{x}^{4}}{81} + \dfrac{{y}^{4}}{16} = 1$ 在曲线上，故曲线 $C$ 关于原点对称，故正确
    \item 取 $x = 0$ ， $y =  \pm  2$ ，取 $y = 0$ ， $x =  \pm  3$ ，故曲线在一个长为 6，宽为 4 的矩形内部，故其面积小于 $6 \times  4 = {24}$ ,故正确
    \item 设曲线上一点为 $M\left( {x,y}\right)$ ，则 $\dfrac{{x}^{4}}{81} + \dfrac{{y}^{4}}{16} = 1$ ，设 $\left\{  \begin{array}{l} {x}^{2} = 9\cos \theta \\  {y}^{2} = 4\sin \theta  \end{array}\right.$
    
    $M$ 到原点的距离的平方为 ${x}^{2} + {y}^{2} = 9\cos \theta  + 4\sin \theta  = \sqrt{97}\sin \left( {\theta  + \varphi }\right) ,\theta  \in  \left\lbrack  {0,\dfrac{\pi }{2}}\right\rbrack, \tan \varphi  = \dfrac{9}{4}$ ，当 $\sin \left( {\theta  + \varphi }\right)  = 1$ 时，距离平方有最大值为 $\sqrt{97}$ ，故距离的最大值为 ${97}^{\dfrac{1}{4}}$ ，故错误
    \item 对于曲线 $C : \dfrac{{x}^{4}}{81} + \dfrac{{y}^{4}}{16} = 1$ 和椭圆 $\dfrac{{x}^{2}}{9} + \dfrac{{y}^{2}}{4} = 1$ ，设点 $\left( {x,{y}_{1}}\right)$ 在 $\dfrac{{x}^{4}}{81} + \dfrac{{y}^{4}}{16} = 1$ 上，

    点 $\left( {x,{y}_{2}}\right)$ 在 $\dfrac{{x}^{2}}{9} + \dfrac{{y}^{2}}{4} = 1$ 上,则 $\dfrac{{x}^{2}}{9} \leq  1$ ,

    ${y}_{1}^{4} - {y}_{2}^{4} = {16}\left( {1 - \dfrac{{x}^{4}}{81}}\right)  - {16}{\left( 1 - \dfrac{{x}^{2}}{9}\right) }^{2} = {16}\left( {1 - \dfrac{{x}^{2}}{9}}\right) \left( {1 + \dfrac{{x}^{2}}{9}}\right)  - {16}{\left( 1 - \dfrac{{x}^{2}}{9}\right) }^{2}$

    $= {16}\left( {1 - \dfrac{{x}^{2}}{9}}\right) \left( {1 + \dfrac{{x}^{2}}{9} - 1 + \dfrac{{x}^{2}}{9}}\right)  = \dfrac{32}{9}{x}^{2}\left( {1 - \dfrac{{x}^{2}}{9}}\right)  \geq  0$ ,故 ${y}_{1}^{4} \geq  {y}_{2}^{4}$ ,所以 $\left| {y}_{1}\right|  \geq  \left| {y}_{2}\right|$ ,

    设点 $\left( {{x}_{1},y}\right)$ 在 $\dfrac{{x}^{4}}{81} + \dfrac{{y}^{4}}{16} = 1$ 上,点 $\left( {{x}_{2},y}\right)$ 在 $\dfrac{{x}^{2}}{9} + \dfrac{{y}^{2}}{4} = 1$ 上,则 $\dfrac{{y}^{2}}{4} \leq  1$ ,

    ${x}_{1}^{4} - {x}_{2}^{4} = {81}\left( {1 - \dfrac{{y}^{4}}{16}}\right)  - {81}{\left( 1 - \dfrac{{y}^{2}}{4}\right) }^{2} = {81}\left( {1 - \dfrac{{y}^{2}}{4}}\right) \left\lbrack  {\left( {1 + \dfrac{{y}^{2}}{4}}\right)  - \left( {1 - \dfrac{{y}^{2}}{4}}\right) }\right\rbrack   = \dfrac{{81}{y}^{2}}{2}\left( {1 - \dfrac{{y}^{2}}{4}}\right)  \geq  0$，所以 ${x}_{1}^{4} \geq  {x}_{2}^{4}$ ,即 $\left| {x}_{1}\right|  \geq  \left| {x}_{2}\right|$

    故椭圆 $\dfrac{{x}^{2}}{9} + \dfrac{{y}^{2}}{4} = 1$ 在曲线 $C : \dfrac{{x}^{4}}{81} + \dfrac{{y}^{4}}{16} = 1$ 内(除四个交点外)，如图:

    \begin{center}
    \includegraphics[width=0.3\textwidth]{bodyxb/src/1_151_874_378_247_0.jpg}
    \end{center}

    设直线 $x - y + \sqrt{5} = 0$ 交椭圆 $\dfrac{{x}^{2}}{9} + \dfrac{{y}^{2}}{4} = 1$ 于 $A,B$ 两点,交 $x$ 轴于 $N\left( {-\sqrt{5},0}\right)$，$M,N$ 为椭圆 $\dfrac{{x}^{2}}{9} + \dfrac{{y}^{2}}{4} = 1$ 的两个焦点，由椭圆的定义可知: $\left| {AN}\right|  + \left| {AM}\right|  = {2a} = 6,\left| {BN}\right|  + \left| {BM}\right|  = {2a} = 6$，所以 $\bigtriangleup  {ABM}$ 的周长为 12，由图可知， $\bigtriangleup  {PQM}$ 的周长大于 12，故正确.
    \end{dingautolist}
    \end{jieda}
\end{enumerate}